% Requires the notation macros of the main manuscript: \E, \Prob, \Law, \Dgm, \norm, \abs.

\section{Setup, estimands, and identification}
\label{sec:setup}

\subsection{Data, sampling unit, and summaries}

The sample consists of $n$ independent point clouds. Unit $i$ contributes the triple
$(X_i, A_i, Y_i)$, $i=1,\dots,n$, an i.i.d.\ draw of $(X,A,Y)$ from a law $P$: the covariates
$X_i \in \mathcal X \subseteq \mathbb R^p$ are measured before the cloud is formed,
$A_i \in \{0,1\}$ is the group indicator with propensity $e(x) = \Prob(A=1 \mid X=x)$, and
$Y_i \in \mathcal Y$ is a point cloud taking values in a Polish space. The sampling unit is the
cloud: every unit contributes exactly one cloud and one summary of it, and no summary, statistic
or calibration draw is formed by pooling points across clouds. Persistent homology is applied to
each cloud separately.

Fix a filtration $F$ and a finite set of homology degrees $\mathcal D$, with
$\mathcal D = \{0,1\}$ the default; the computational study uses an
alpha-complex filtration computed separately on each cloud.
Write $D_{i,d} = \Dgm_d(F(Y_i))$ for the degree-$d$ diagram,
a finite multiset of pairs $p=(a_p,b_p)$ with $0 \le a_p < b_p < \infty$ after essential classes
are dropped. Diagrams are computed once per cloud and cached, so no persistent homology is
recomputed inside a calibration loop. For a fixed compact interval $T \subset \mathbb R_+$ and a
weight power $r>0$, the degree-$d$ power-weighted silhouette is
\begin{equation}
\label{eq:silhouette}
\begin{split}
Z_{i,d}(t) = \Phi_d(D_{i,d})(t)
&= \frac{\sum_{p \in D_{i,d}} w_p \Lambda_p(t)}{\sum_{p \in D_{i,d}} w_p},
\qquad
w_p = (b_p - a_p)^r,
\\
\Lambda_p(t) &= \max\{0, \min(t - a_p,\, b_p - t)\},
\end{split}
\end{equation}
evaluated on a fixed grid $t \in T$; an empty diagram contributes the zero function. The weight
power is $r=3$ throughout, and the silhouette grid is fixed before any group contrast. As
vector-valued summaries of the per-cloud diagram law, silhouettes obey a law of large numbers and
a central limit theorem under integrability and stability conditions
\cite{bubenik_jmlr_2015,chazal_socg_2014,kim_lee_tce_2026}, which is why they serve as the
outcome representation below.

\subsection{Potential outcomes and the covariate-standardized targets}

Let $Y_i^a$ be the potential cloud under assignment $a \in \{0,1\}$, with potential diagrams
$D_{i,d}^a = \Dgm_d(F(Y_i^a))$ and potential silhouettes $Z_{i,d}^a(t) = \Phi_d(D_{i,d}^a)(t)$.

\begin{assumption}[Identification]
\label{ass:identification}
\emph{(i) Causal consistency:} $Y_i = Y_i^{A_i}$, so
$D_{i,d} = D_{i,d}^a$ and $Z_{i,d}(t) = Z_{i,d}^a(t)$ on $\{A_i=a\}$. \emph{(ii) Arm-specific
conditional exchangeability:} $Y_i^a \perp A_i \mid X_i$ for each $a \in \{0,1\}$. \emph{(iii)
Strict positivity:} $0 < e(x) < 1$ for $P_X$-almost every $x$. In addition,
$\E\norm{Z_{i,d}^a}_\infty < \infty$ for each $a$ and $d$.
\end{assumption}

Define the arm-specific outcome regression $m_{a,d}(t,x) = \E[Z_{i,d}^a(t) \mid X_i = x]$ and the
covariate-standardized topological effect
\begin{equation}
\label{eq:psi}
\psi_d(t) = \E_X\bigl[m_{1,d}(t,X) - m_{0,d}(t,X)\bigr]
= \E[Z_{i,d}^1(t)] - \E[Z_{i,d}^0(t)].
\end{equation}
Under Assumption~\ref{ass:identification} the g-formula gives $\E[Z_{i,d}^a(t)] = \E[m_{a,d}(t,X)]$, and the inverse
propensity representation $\E[Z_{i,d}(t)\mathbb 1\{A_i=a\}/e_a(X_i)]$ equals the same quantity,
with $e_1=e$ and $e_0=1-e$ \cite{souto_diamantis_tcda_2026,kim_lee_tce_2026}. Hence $\psi_d(t)$ is
identified from the observational law. Only the arm-specific form of exchangeability is used;
joint exchangeability of $(Y^0,Y^1)$ would be needed only for joint targets that this paper does
not estimate.

At the distribution level, let $\mu_D = \sum_{p \in D} w_p \delta_{\mathrm{coord}(p)}$ be the
weighted persistence measure of a diagram, with $\mathrm{coord}(p)$ the chosen diagram
coordinates, and let $T_{\mathrm{dist}}(P) = \E_P[\mu_{D(F(Y))}]$ be the expected persistence
measure \cite{divol_chazal_2020,divol_lacombe_2021}. With $Q_a(x,\cdot) = \Law(Y \mid A=a, X=x)$,
the law-valued g-formula $P^a_Y = \int_{\mathcal X} Q_a(x,\cdot)\,dP_X(x)$ holds under the same
Assumption~\ref{ass:identification}, so $T_{\mathrm{dist}}(P^a_Y) = \int_{\mathcal X}
\E[\mu_{D(F(Y))} \mid A=a, X=x]\,dP_X(x)$ \cite{souto_diamantis_tcda_2026}. Two facts about this
construction are load-bearing. First, the measure is un-normalized and therefore affine in $P$,
so the covariate mixture commutes with it; the normalized silhouette of \eqref{eq:silhouette} is
not affine, and persistent homology is in any case applied per cloud before the mixture. Second,
the expected measure is evaluated through its dual pairing with a fixed separable test class
rather than as a Bochner integral in total variation, which is not separable; the fixed grid used
below is exactly such a finite separable class.

The implemented distribution-level target is the projection of $T_{\mathrm{dist}}$ onto one frozen
grid. The grid has $32 \times 32$ bins on the interval $(0,2)$ in the (birth, midpoint) plane,
with weight power $r=3$, giving $q = 1024$ coordinates. Let $\Pi_G$ denote the corresponding
binning map and
\begin{equation}
\label{eq:delta}
\delta = \Pi_G T_{\mathrm{dist}}(P^1_Y) - \Pi_G T_{\mathrm{dist}}(P^0_Y) \in \mathbb R^q.
\end{equation}
The shipped tests concern $H_0^{\mathrm{dist,grid}} : \delta = 0$ only. Equality of the
unprojected measures is not certified without a resolution or injectivity argument for $\Pi_G$,
and no such claim is made. The same assumption set identifies both levels, with no exchangeability
or positivity asymmetry between the outcome-level and distribution-level targets. If a reader
replaces arm-specific conditional exchangeability by conditional topological ignorability alone,
the unstandardized interventional contrast is not identified and the identified object drops to a
covariate-standardized within-stratum contrast \cite{saki_faghihi_2026}; under Assumption~\ref{ass:identification} as
stated, both covariate-standardized targets above are identified.

\subsection{Null hypotheses and their relations}

The paper's tests are indexed by three population nulls, together with one unprojected reference statement. The field's null is the conditional-law null
\begin{equation}
H_0^{\mathrm{cond}} : \Law(D_{i,d} \mid A_i=1) = \Law(D_{i,d} \mid A_i=0)
\quad \text{for every } d \in \mathcal D,
\end{equation}
which is testable without a causal model and is the target of the existing topological two-sample
tests
\cite{robinson_turner_2017,kwitt_neurips_2015,han_kim_kim_2026,murris_strand_2026,moon_lazar_2023,dubey_muller_2019}.
The covariate-standardized outcome null and the finite-grid distribution null are
\begin{equation}
H_0^{\mathrm{out}} : \psi_d(t) = 0 \ \text{ for all } t \in T \text{ and all } d \in \mathcal D,
\qquad
H_0^{\mathrm{dist,grid}} : \delta = 0.
\end{equation}
We also write $H_0^{\mathrm{dist}}$ for the unprojected distribution-level statement. Two versions of it appear below: the expected-measure version, equality of $T_{\mathrm{dist}}(P^1_Y)$ and $T_{\mathrm{dist}}(P^0_Y)$, which enters the pairing relation to $H_0^{\mathrm{out}}$ and the exact witnesses, and the law-equality version, $\Law(Y^1)=\Law(Y^0)$, which the characteristic-kernel MMD pairing targets. Neither unprojected statement is certified by the finite-grid test.

The relation between the unprojected $H_0^{\mathrm{dist}}$ and $H_0^{\mathrm{out}}$ is pairing-conditional, not a fact about topology. If the
outcome-level mean factors through the distribution-level representation, that is, if the mean
silhouette functional equals $L \circ T_{\mathrm{dist}}$ for some map $L$ and the metric separates
points, then $H_0^{\mathrm{dist}} \Rightarrow H_0^{\mathrm{out}}$. For the default pairing of this
paper, the normalized power-weighted silhouette against the expected persistence measure, no such
factorization exists, because the normalization ratio sits inside the expectation, and neither
null implies the other. Two exact witnesses establish the logical independence. In W1,
$D^0 = \{(0,\ell)\}$ and $D^1 = \{(0,\ell),(0,\ell)\}$ for a fixed $\ell>0$: one homology class
splits into two classes at the same merge scale. The normalized silhouette averages the tents, so
the mean silhouettes are equal and $H_0^{\mathrm{out}}$ holds, while the expected measures carry
masses $w_\ell$ and $2w_\ell$ at the same coordinate and are unequal, so
$H_0^{\mathrm{dist,grid}}$ fails; on the frozen grid the projected contrast is $2.460375$ while
the mean-silhouette contrast is exactly $0$. In W2$'$, $D^0 = \{(0,1),(0,2)\}$ and $D^1$ equals
$\{(0,1),(0,1)\}$ or $\{(0,2),(0,2)\}$ with probability $1/2$ each: each location carries expected
multiplicity one in both arms, so the expected measures agree exactly for every weight power and
$H_0^{\mathrm{dist,grid}}$ holds, while the two arms weight the two tents by
$(w_1/(w_1+w_2),\, w_2/(w_1+w_2))$ against $(1/2,1/2)$, with $w_j = j^r$, so the mean silhouettes
differ and $H_0^{\mathrm{out}}$ fails; the projected measure contrast is exactly $0$ and the
mean-silhouette contrast is $0.544$. The mechanism in both directions is the non-commutation of
the normalization ratio with mixing. W1 is a statement about the outcome-level mean only: it does
not say that an unadjusted test of $H_0^{\mathrm{cond}}$ is blind, since the arm-conditional laws
differ there.

A secondary distribution-level variant replaces the mean by the squared maximum mean discrepancy
under the exponentiated Reininghaus persistence scale-space kernel
$k^U_\sigma = \exp(k_\sigma)$. On the set of diagrams whose birth and death coordinates are
bounded by a fixed constant and whose total multiplicity is bounded by a fixed constant, this
kernel is universal with respect to the Wasserstein-1 metric \cite{kwitt_neurips_2015};
universality implies characteristic, so its mean embedding is injective on laws and its null is
equality of the interventional diagram laws \cite{gretton_2012}. That null is strictly stronger
than $H_0^{\mathrm{out}}$, in the sense that it implies $H_0^{\mathrm{out}}$ with a false converse.
It is reported as a secondary with its nested target flagged, both because it checks whether an
effect was integrated away by the mean and because its ordering against the outcome null differs
from the headline pairing; W1 and W2$'$ separate under it with squared discrepancies $1.916989$
and $0.398762$.

Comparisons of rejection rates across these nulls are not efficiency statements, and the paper
reports none. Every performance comparison is made within a fixed target, and the failure of the
unadjusted tests is a target-mismatch result, stated in the theory and experiments sections,
rather than a power ranking.

\section{Inference methods}
\label{sec:methods}

\subsection{Outcome-level test}

Let $\hat e$ and $\hat m_{a,d}$ denote cross-fitted nuisance estimates, each unit's prediction
coming from folds that exclude it, with folds stratified in $A$. The per-unit cross-fitted augmented inverse-probability-weighted (AIPW)
score in degree $d$ is
\begin{equation}
\label{eq:outscore}
s_{i,d}(t) = \hat m_{1,d}(t,X_i) - \hat m_{0,d}(t,X_i)
+ \frac{A_i}{\hat e(X_i)}\bigl\{Z_{i,d}(t) - \hat m_{1,d}(t,X_i)\bigr\}
- \frac{1-A_i}{1-\hat e(X_i)}\bigl\{Z_{i,d}(t) - \hat m_{0,d}(t,X_i)\bigr\},
\end{equation}
with sample mean $\hat\psi_d(t) = n^{-1}\sum_i s_{i,d}(t)$. The score is the efficient-influence
function linearization of the contrast and is doubly robust: its expectation is $\psi_d(t)$ if
either the propensity or the outcome regression is correct, and the remainder is a product of
nuisance errors \cite{souto_diamantis_tcda_2026,chernozhukov_dml_2018,kennedy_drlearner_2023}. The
point estimator, the cross-fitting and the per-unit influence process are reused from the
released implementation accompanying the framework of \cite{souto_diamantis_tcda_2026}; they are
not reimplemented here, and this paper owns the testing layer, namely the statistic, the null law
and the p-value.
The statistic is
\begin{equation}
\label{eq:tnout}
T_n^{\mathrm{out}} = \sqrt n \max_{d \in \mathcal D} \sup_{t \in T} \abs{\hat\psi_d(t)}.
\end{equation}

The null law is a Gaussian multiplier bootstrap over the centered score process. Draw one
multiplier $\xi_{i,b} \sim N(0,1)$ per unit, shared across all degrees, and set
\[
G_b = \max_{d \in \mathcal D} \sup_{t \in T} \abs{n^{-1/2}\sum_i \xi_{i,b}\bigl(s_{i,d}(t) - \bar
s_d(t)\bigr)}.
\]
The shared streams preserve the dependence across degrees induced by the shared units
and folds; drawing independent multipliers per degree would replace that dependence with
independence and inflate the null. The Gaussian multiplier approximation for maxima of sums is the
one of \cite{cck_2013}, and the p-value applies the Phipson-Smyth correction,
$(1 + \#\{b : G_b \ge T_n^{\mathrm{out}}\})/(1+B)$. The default implementation uses
$B=2000$ draws; the reported oracle and cloud batteries use $399$ multiplier draws
with $199$ frozen stratified permutations, and the distribution-level test uses
$1999$ draws (Section~\ref{sec:exp:design}) \cite{phipson_smyth_2010}.

As an alternative calibration in an exchangeability regime, the same statistic can be referred to
the frozen stratified permutation null: labels are permuted within propensity strata while the
fitted nuisances, the strata and the propensity clipping are held fixed, and the score is
re-evaluated under the permuted labels without refitting. This is conditionally exact
randomisation inference when the nuisances and strata are fixed independently of the permuted
labels, for instance under a sharp conditional null with externally fixed design information; with
estimated, label-dependent nuisances it is a fast cross-fitted calibration approximation, not a
finite-sample theorem \cite{chung_romano_2013,janssen_1997}. The paper reports it as an option and
does not claim exactness in the general estimated-nuisance regime.

\subsection{Distribution-level test}

Let $V_i \in \mathbb R^q$ collect the binned weighted masses of the degree-$d$ persistence measure
on the frozen grid of \eqref{eq:delta}. Let $\hat m_a$ be the cross-fitted regression of $V$ on
$X$ within arm $a$, and let $\hat e$ be the cross-fitted propensity. The per-unit AIPW contrast
score on the fixed grid is
\begin{equation}
\label{eq:distscores}
\varphi_i = \hat m_1(X_i) - \hat m_0(X_i)
+ \frac{A_i}{\hat e(X_i)}\bigl\{V_i - \hat m_1(X_i)\bigr\}
- \frac{1-A_i}{1-\hat e(X_i)}\bigl\{V_i - \hat m_0(X_i)\bigr\},
\end{equation}
with sample mean $\hat\delta = n^{-1}\sum_i \varphi_i$. Its coordinates are the finite-dimensional
specialization of the Banach-valued doubly robust representation, and the same delegation
boundary as at the outcome level applies: the nuisances and the influence linearization are
reused, while the statistic, null law and p-value are this paper's. Coordinatewise variances use
the $1/n$ normalization,
$\hat\sigma_j^2 = n^{-1}\sum_i (\varphi_{ij} - \hat\delta_j)^2$, and the statistic is the
coordinatewise studentized maximum
\begin{equation}
\label{eq:tndist}
T_n^{\mathrm{dist}} = \max_{j \in J} \frac{\sqrt n \, \abs{\hat\delta_j}}{\hat\sigma_j},
\qquad
J = \{j : \hat\sigma_j^2 > 10^{-8}\}.
\end{equation}
Coordinates with sample variance at or below the floor $10^{-8}$ are dropped from the statistic
and from the null, and a sample in which every coordinate is dropped raises a deterministic
failure rather than passing silently. The null uses $B=1999$ Gaussian multiplier draws, with one
multiplier per unit shared across all coordinates: draw $b$ is
$\max_{j \in J} \abs{n^{-1/2}\sum_i \xi_{i,b}(\varphi_{ij} - \hat\delta_j)}/\hat\sigma_j$, and the
p-value again uses the Phipson-Smyth correction \cite{cck_2013,phipson_smyth_2010}. No persistent
homology or nuisance regression is recomputed inside the calibration loop. As a basic check, with
balanced arms, a known propensity of one half and in-sample arm-mean nuisances, the statistic
reduces exactly to the ordinary pooled two-sample $t$ statistic on a single active coordinate, up
to the factor $\sqrt{n/(n-2)}$.

The scope of this test is exactly $H_0^{\mathrm{dist,grid}} : \delta = 0$, the weak
expected-measure null, which permits unequal unit-level laws; no sharp-null or law-equality claim
is attached to it. The target is the projection onto the single frozen grid, so equality of the
unprojected interventional measures is not certified. In the calibration study the strongest
tested confounding produces one mildly anti-conservative cell at the smallest sample size
($0.096$ at $n=200$), while the same design is inside the nominal band at $n=500$ ($0.066$); the
conservative operating recommendation is therefore $n \ge 500$ under strong confounding, and the
paper claims no uniform control at smaller sample sizes. Permutation calibration for this null is
retained only as a diagnostic: under the weak null with imbalance the permuted score mean is
centered at the observational contrast rather than at zero, so the frozen relabeling law is not
the weak-null sampling law.

\subsection{Multiplicity across homology degrees}

All degree-level procedures read the same null matrix, so differences among them are differences
of procedure rather than of calibration draws. For each null draw $b$, one shared multiplier
vector or one shared label permutation produces a per-degree statistic $X_{d,b}$; the observed
per-degree statistics are $X_d = \sqrt n \sup_t \abs{\hat\psi_d(t)}$. Three reductions are
reported. The unstudentized shared maximum uses $T = \max_d X_d$ with null $\max_d X_{d,b}$; it
preserves cross-degree dependence but lets the loudest degree dominate the family. Bonferroni
multiplies the smallest per-degree rank p-value by $\abs{\mathcal D}$ and is the transparent
conservative comparator. The adapted studentized comparator standardizes each degree against its
own empirical null and then takes the joint maximum; the standardization uses the pooled
convention, with location and scale constants computed from the observed statistic together with
the $B$ null replicates, because only that symmetric convention keeps the rank test exact under
exchangeability. A pairwise Kolmogorov-Smirnov statistic between standardized per-degree null
samples is reported as a descriptive comparability measure only, since the draws are dependent
across degrees and a p-value from it would not be valid.

The studentized comparator is an adaptation of the empirical-null multiple-testing construction of
\cite{vejdemo_mukherjee_2022}, not the original procedure. Three elements are replaced. The source
null, a uniform point process on an estimated convex body, has no counterpart for a
covariate-adjusted treatment-effect null; the source simulates each family member independently,
whereas here the degrees are drawn jointly from one shared multiplier or permutation stream; and
the source standardizes with null-only constants, whereas here the constants are pooled over the
observed statistic and the null replicates. The validity of the adaptation rests on exchangeability
of the null mechanism together with the functional weak convergence of the AIPW process
\cite{kim_lee_tce_2026}, not on the source's point-process limit theorem. No finite-sample
exactness is claimed for the studentized comparator; the unstudentized maximum and Bonferroni
lines are not procedures of the source, and the reported calibration record, including its
conservative misses, is stated with the experiments.
